1つ面白い話をしよう.
まず,次の3つの漸化式を見てほしい.
\[
  a_{n+1}=a_n^2,
  \quad
  a_{n+1}=a_n^2-1,
  \quad
  a_{n+1}=a_n^2-2.
\]
1つ目はきちんとここまで本書を読んできた諸君らならば,用意に解くことができるだろう.
実際,これは対数型漸化式であるから,両辺の対数を取って等比数列に帰着させることができる.

2つ目はどうだろうか$a_n$に指数がついているから,対数を使って処理したいが$-1$の項が邪魔していてうまくいかない.
不動点を見つけようにも指数の項が邪魔をするし,うまくいきそうにない.
少々意地悪であったかもしれないが,これは初等的な解法を用いては,その一般項を得られない漸化式である.

3つ目も一見先の例のように一般項が得難い見た目をしている.
しかしながらこれはある特殊な条件の噛み合いによって偶然にも初等的な解法が存在する.

ここまで,等差型,等比型,階差型,階比型など,漸化式を一定の型に
帰着して解く方法を学んできた.しかし,右辺が少し複雑になると,
同じ道具が急に働かなくなる.

ここで扱うのは,一般の漸化式がいつでも解けるという話ではない.
一見すると解けそうにない式の中にの数学的対象と偶然一致するものがある.
その一致を見抜くことで,初等的な型に還元できない漸化式を解いていく.
本章はそういう話である.